% ============================================================
%  Show Your Work
%  Building a verified math notebook in Lean 4
%  — the book written from building mathbook (github.com/adekau/mathbook)
%  Build: xelatex -shell-escape show-your-work.tex (twice), or latexmk -xelatex -shell-escape
% ============================================================
\documentclass[12pt, a4paper, openright]{book}

% ── Fonts & encoding ────────────────────────────────────────
\usepackage{fontspec}
\IfFontExistsTF{Latin Modern Roman}{\setmainfont{Latin Modern Roman}}{}
% DejaVu Sans Mono has every glyph the Lean listings need (⊔ ⊓ λ ⋖ ≤ ⁻¹); fall back to Latin Modern Mono.
\IfFontExistsTF{DejaVu Sans Mono}
  {\setmonofont[Scale=0.82]{DejaVu Sans Mono}}
  {\IfFontExistsTF{Latin Modern Mono}{\setmonofont[Scale=0.9]{Latin Modern Mono}}{}}

% ── Layout ──────────────────────────────────────────────────
\usepackage[top=1in, bottom=1.15in, left=1.1in, right=1.1in, headheight=14pt]{geometry}
\usepackage{microtype}
\usepackage{placeins}
\usepackage{setspace}
\setstretch{1.08}
\usepackage{parskip}
\setlength{\parskip}{6pt}

% ── Math ────────────────────────────────────────────────────
\usepackage{amsmath}
\usepackage{amssymb}
\usepackage{mathtools}
\usepackage{unicode-math}
\IfFontExistsTF{Latin Modern Math}{\setmathfont{Latin Modern Math}}{}

% ── Theorems ────────────────────────────────────────────────
\usepackage{amsthm}
\newtheoremstyle{defstyle}{8pt}{6pt}{\normalfont}{0pt}{\bfseries}{.}{.5em}{}
\newtheoremstyle{thmstyle}{8pt}{6pt}{\itshape}{0pt}{\bfseries}{.}{.5em}{}
\theoremstyle{thmstyle}
\newtheorem{theorem}{Theorem}[chapter]
\newtheorem{lemma}[theorem]{Lemma}
\newtheorem{corollary}[theorem]{Corollary}
\newtheorem{proposition}[theorem]{Proposition}
\theoremstyle{defstyle}
\newtheorem{definition}[theorem]{Definition}
\newtheorem{example}[theorem]{Example}
\newtheorem{remark}[theorem]{Remark}

% ── Code (minted) ───────────────────────────────────────────
\usepackage{minted}
\usemintedstyle{friendly}
\setminted[lean]{
  bgcolor=gray!6, frame=lines, framesep=4pt, rulecolor=gray!40,
  fontsize=\small, breaklines=true, linenos=false, xleftmargin=8pt,
}
\setminted[text]{
  bgcolor=yellow!8, frame=single, framesep=4pt, rulecolor=gray!30,
  fontsize=\small, linenos=false, xleftmargin=6pt, breaklines=true,
}
\setminted[typescript]{
  bgcolor=gray!6, frame=lines, framesep=4pt, rulecolor=gray!40,
  fontsize=\small, breaklines=true, linenos=false, xleftmargin=8pt,
}
\newcommand{\lc}[1]{\mintinline{lean}{#1}}

% ── Diagrams ────────────────────────────────────────────────
\usepackage{tikz}
\usetikzlibrary{positioning, arrows.meta, shapes.geometric, calc, decorations.pathreplacing}
\tikzset{
  node/.style={circle, draw=gray!80, fill=gray!10, inner sep=3pt, minimum size=7mm, font=\small},
  edge/.style={-{Stealth[scale=0.8]}, gray!60, thick},
  hasseedge/.style={draw=gray!70, thick},
  box/.style={draw=gray!70, rounded corners=3pt, fill=gray!5, inner sep=6pt, font=\small, align=center},
}

% ── Coloured boxes ──────────────────────────────────────────
\usepackage{tcolorbox}
\tcbuselibrary{theorems, breakable, skins, most}
\tcbset{
  defbox/.style={enhanced, breakable, colback=blue!4, colframe=blue!55!black, fonttitle=\bfseries,
    boxrule=0.6pt, arc=3pt, left=6pt, right=6pt, top=6pt, bottom=6pt},
  thmbox/.style={enhanced, breakable, colback=green!3, colframe=green!50!black, fonttitle=\bfseries,
    boxrule=0.6pt, arc=3pt, left=6pt, right=6pt, top=6pt, bottom=6pt},
  exbox/.style={enhanced, breakable, colback=orange!4, colframe=orange!70!black, fonttitle=\bfseries,
    boxrule=0.6pt, arc=3pt, left=6pt, right=6pt, top=6pt, bottom=6pt},
  exercisebox/.style={enhanced, breakable, colback=purple!4, colframe=purple!60!black, fonttitle=\bfseries,
    boxrule=0.6pt, arc=3pt, left=6pt, right=6pt, top=6pt, bottom=6pt},
  notebox/.style={enhanced, colback=yellow!8, colframe=yellow!70!black,
    boxrule=0.4pt, arc=3pt, left=6pt, right=6pt, top=4pt, bottom=4pt},
  logbox/.style={enhanced, breakable, colback=gray!6, colframe=gray!60, fonttitle=\bfseries,
    boxrule=0.5pt, arc=3pt, left=6pt, right=6pt, top=6pt, bottom=6pt},
  trybox/.style={enhanced, breakable, colback=teal!4, colframe=teal!50!black, fonttitle=\bfseries,
    boxrule=0.6pt, arc=3pt, left=6pt, right=6pt, top=6pt, bottom=6pt},
}
\newenvironment{defbox}[1][]{\begin{tcolorbox}[defbox, title={Definition\ifx\\#1\\\else\ --- #1\fi}]}{\end{tcolorbox}}
\newenvironment{thmbox}[1][]{\begin{tcolorbox}[thmbox, title={\ifx\\#1\\Theorem\else#1\fi}]}{\end{tcolorbox}}
\newenvironment{exbox}[1][]{\begin{tcolorbox}[exbox, title={Example\ifx\\#1\\\else\ --- #1\fi}]}{\end{tcolorbox}}
\newenvironment{notebox}{\begin{tcolorbox}[notebox]}{\end{tcolorbox}}
% What actually happened while building — the entries of book/TRACKING.md, retold.
\newenvironment{logbox}[1][]{\begin{tcolorbox}[logbox, title={From the log\ifx\\#1\\\else\ --- #1\fi}]}{\end{tcolorbox}}
% Cells to type into the notebook.
\newenvironment{trybox}[1][]{\begin{tcolorbox}[trybox, title={Try it\ifx\\#1\\\else\ --- #1\fi}]}{\end{tcolorbox}}
\newcounter{exercise}[chapter]
\newenvironment{exercisebox}[1][]{\refstepcounter{exercise}\begin{tcolorbox}[exercisebox,
    title={Exercise \thechapter.\theexercise\ifx\\#1\\\else\ --- #1\fi}]}{\end{tcolorbox}}

% ── Chapter headings, headers, links ────────────────────────
\usepackage{titlesec}
\titleformat{\chapter}[display]{\huge\bfseries}{\large\scshape\chaptertitlename\ \thechapter}{10pt}{\Huge\bfseries}
\titlespacing*{\chapter}{0pt}{20pt}{30pt}
\usepackage{fancyhdr}
\pagestyle{fancy}
\fancyhf{}
\fancyhead[LE]{\small\itshape\leftmark}
\fancyhead[RO]{\small\itshape\rightmark}
\fancyfoot[C]{\small\thepage}
\renewcommand{\headrulewidth}{0.4pt}
\usepackage[colorlinks=true, linkcolor=blue!60!black, citecolor=green!50!black, urlcolor=blue!70!black,
            bookmarks=true, bookmarksnumbered]{hyperref}
\usepackage{enumitem}
\usepackage{booktabs}
\usepackage{longtable}
\usepackage{emptypage}
\setlength{\emergencystretch}{2em}
\usepackage{xspace}
\newcommand{\rn}[1]{\texttt{#1}}
\newcommand{\file}[1]{\path{#1}}
\newcommand{\status}[1]{\textsc{#1}}

% ============================================================
\begin{document}
\frontmatter

\begin{titlepage}
  \centering
  \vspace*{2cm}
  {\Huge\bfseries Show Your Work\par}
  \vspace{0.5cm}
  {\large\itshape Building a verified math notebook in Lean~4\par}
  \vspace{3cm}
  \begin{tcolorbox}[width=0.75\textwidth, colback=gray!8, colframe=gray!50, arc=4pt, boxrule=0.6pt]
    \small\centering
    \textit{Every line a computer-algebra system prints is a claim.\\
    This book is about a notebook that shows its work and says,\\
    for each step, exactly how much of the claim is a theorem.}
  \end{tcolorbox}
  \vfill
  {\large One engine, in Lean~4, compiled to native and to WebAssembly,\\
  and the book written from building it\par}
  \vspace{1cm}
  {\normalsize The engine imports nothing but Lean's prelude; the real-number theorems live beside it, in Mathlib.\par}
  \vspace{0.5cm}
  {\small\color{gray} Edition 1.0 · mathbook at commit \texttt{\input{chapters/commit.tex}}\par}
\end{titlepage}

\chapter*{Preface}
\addcontentsline{toc}{chapter}{Preface}

This book was written from building a thing: a notebook, in the spirit of Mathematica, that
differentiates, integrates, row-reduces and simplifies, and that \emph{shows its work} --- every
answer arrives with the chain of rewrites that produced it, every subterm of every line can be
clicked to ask where it came from, and every rule that fired carries a label saying how much of
it is a theorem. The engine behind the notebook is a single Lean~4 program. It runs natively as a
server and a command-line tool, and it runs in the browser as WebAssembly; the notebook never
computes anything itself. The proofs about the engine are Lean~4 too, and they are checked every
time the project builds.

The order of the chapters is roughly the order in which the thing was built, because the order
turned out to matter. A rewriter came first, and with it a trace. Then meaning: what each rule
claims, over the integers, then over the reals, then under a derivative. Then termination,
which was the hard part and is the centre of the book: the engine has no step budget, and the
proof that it needs none took several wrong orderings and one wall that would not move. Then the
tracing of origins, then three more ``worlds'' --- linear algebra, integration, the
$\lambda$-calculus and finite order theory --- each of which had to find its own honest place
between \emph{proved} and \emph{checked}.

\medskip
\textbf{Who this book is for.} Someone who can read undergraduate mathematics and some functional
code, who has perhaps met Lean~4 in one of its companion volumes, and who wants to see what
``verified'' can mean for a system that is also meant to be \emph{used}. No Mathlib is needed to
read the engine; where the reals enter, the book says so.

\medskip
\textbf{What is claimed.} Precisely what the notebook's status labels claim, and no more. A rule is
\status{verified} when an unconditional theorem says its output means what its input meant;
\status{conditional} when the theorem needs a hypothesis and the \emph{necessity} of the hypothesis
is itself proved; \status{checked} when the step is a guess that a later step verifies; and
\status{unverified} when there is no theorem yet. Chapter~\ref{ch:ledger} is the full ledger.
Several places where a proof was avoided are pointed at by name: they are the run-time checks the
book calls \emph{honest boundaries}, and the reader will see exactly what is checked and why a
check was the right thing there.

\medskip
\textbf{How to read this book.} Part~I is the engine's shape and the rewriter that keeps a trace.
Part~II gives the rules their meaning. Part~III is termination. Part~IV is the notebook side of
``show your work''. Part~V adds the other worlds. Boxes labelled \emph{From the log} retell what
actually happened --- the wrong turns are kept, since they are the argument for the right ones ---
and boxes labelled \emph{Try it} give cells to type into the notebook.

\medskip
\textbf{The code.} Everything quoted is from the repository at the commit stamped on the title
page. File names are given the way the repository has them: \file{engine/MathEngine/Order.lean}
is the ordering, \file{proofs/Proofs/SimpReal.lean} the real-number theorems, and so on.

\tableofcontents

\mainmatter
\part{One Engine}
\chapter{What a Verified Notebook Is}\label{ch:what}

\begin{quote}\itshape
``The answer is $2x\cos x + x^2\sin x$.''\;---\;``Says who?''
\end{quote}

A computer-algebra system prints an answer and asks to be believed. Usually it is right, and
usually that is enough. This book is about a notebook built on a different bargain: it prints the
answer \emph{and} the steps, and for each step it says what would have to be true for the step
to be a lie. Sometimes nothing: the step is a theorem. Sometimes one named hypothesis. Sometimes
``this step is a guess, and the step below it checks the guess''. And sometimes ``no theorem yet''.
The notebook never says more than it can back.

\section{Show your work}

Type \texttt{diff(x\^{}2 * sin(x), x)} into the notebook and it answers
$2x\sin x + x^2\cos x$, with seven steps underneath: the product rule, the power rule, a constant
folded, an identity applied, the chain rule, an identity applied again. Each step is a line, the
whole term after the step, with the part that changed underlined. Click any subterm of any line
--- the $\cos x$ in the answer, say --- and a panel names the step that \emph{created} it (the
chain rule), the steps that merely \emph{copied} it along, and the steps that fired \emph{inside}
it, together with the proof status of each. That panel is Chapter~\ref{ch:origins}.

The four statuses are the book's spine, so they are worth stating once, properly.

\begin{defbox}[The four statuses]
  Every rule the engine can fire is reported to the notebook with one of four labels.
  \begin{description}[nosep]
    \item[\status{verified}] An unconditional theorem states that the output of the rule means the
      same as its input, in the semantics of the chapter that proved it (the integers, the reals,
      the reals under a derivative, the rationals in a linear system, \dots).
    \item[\status{conditional}] The theorem needs a side condition --- \emph{and the side condition
      is shown to be necessary}: a counterexample is proved, not just suspected. The notebook shows
      the condition.
    \item[\status{checked}] The step is a guess. A later step verifies the guess by an independent
      computation whose own steps carry their own statuses. (Integration, Chapter~\ref{ch:integrate}.)
    \item[\status{unverified}] There is no theorem yet. The notebook says so, in those words.
  \end{description}
  A step with nested steps inherits the weakest status below it.
\end{defbox}

The distinction between the first two is the one that gives the book most of its stories. A
computer-algebra system simplifies $x/x$ to $1$. So does this engine. But $x/x$ is not $1$ at
$x = 0$, and rather than pretend otherwise, or refuse the simplification everyone expects, the
engine keeps the rule and \emph{writes down the assumption}: the rule is \status{conditional} on
the merged base being nonzero, and a theorem in Chapter~\ref{ch:sound} proves that without the
assumption the rule turns $0$ into $1$. That is the whole attitude of the project in one rule.

\section{One engine}

There is exactly one implementation of the mathematics, and it is written in Lean~4. It compiles
to a native executable (behind an HTTP host, or as a command-line tool speaking JSON-RPC on
standard input) and, from the same source, to WebAssembly running in a web worker. The notebook
is a static page; it draws terms with KaTeX and draws plots and Hasse diagrams as SVG, and it never
evaluates anything. Every number, every rewrite, every ``this is a lattice'' comes over the wire
from the engine, with the derivation attached.

\begin{center}
\resizebox{\linewidth}{!}{\begin{tikzpicture}[node distance=9mm]
  \node[box] (nb) {notebook (static page)\\KaTeX, SVG, no mathematics};
  \node[box, right=9mm of nb] (proto) {JSON-RPC over a Transport\\worker · HTTP · stdio};
  \node[box, right=9mm of proto] (eng) {the engine, Lean 4\\native and wasm};
  \node[box, below=of eng] (proofs) {\file{proofs/}: theorems about\\the engine, with Mathlib};
  \draw[edge] (nb) -- (proto);
  \draw[edge] (proto) -- (eng);
  \draw[edge, dashed] (proofs) -- node[right, font=\scriptsize] {imports} (eng);
\end{tikzpicture}}
\end{center}

Two decisions from the first week held for the whole project.

\paragraph{Two packages.} The engine, in \file{engine/}, imports nothing but Lean's prelude
\lc{Init}. The reason is mechanical: the WebAssembly build emits C for every imported Lean module,
and \lc{Init} alone is 631 of them; Mathlib would be tens of thousands. So the real numbers,
which need Mathlib, cannot live in the engine. They live in a second package, \file{proofs/},
which imports the engine and Mathlib and contains theorems only. A script,
\file{scripts/check-engine-deps.sh}, fails the build if the engine ever acquires a dependency.
The consequence for the reader: every theorem in this book either lives in the engine (and is
proved with Lean's core tactics --- \lc{omega}, \lc{simp}, \lc{grind}, \lc{decide}) or lives in
\file{proofs/} and may use anything Mathlib knows about $\mathbb{R}$. The book says which.

\paragraph{The protocol is the seam.} The notebook and the engine share a small JSON-RPC contract
(\file{packages/protocol}). Its one design rule --- rule~5 --- is that new capabilities are added
as \emph{optional fields and optional methods}, never as changes. Every chapter of this book that
added something to the wire (paths, proof statuses, origin traces, the de~Bruijn view, the Hasse
diagram) added it as a field a frontend may ignore. The notebook from week one would still run
against the engine from the last chapter.

\section{The spike: does the engine reach the browser?}

Before any mathematics, one question decided whether the project was possible: can a Lean~4
program of this size run in a browser, on \emph{current} Lean, fast enough to answer as you type?

\begin{logbox}[Milestone 0]
  Lean stopped shipping a prebuilt 32-bit WebAssembly runtime after v4.15. The project's policy is
  the latest stable toolchain (v4.33.1 at the time of writing), so the runtime and \lc{Init} were
  built from source for \texttt{wasm32} with Emscripten: about nine minutes the first time, cached
  afterwards. The result was a 1.4~MB \texttt{engine-lean.wasm} that answered its first request
  96~ms after the worker started, and evaluated \texttt{x + 0} in 50~µs per call after warm-up. No
  threads, so no cross-origin isolation headers: the notebook stays a plain static directory.

  One upstream bug surfaced on the way. On a 32-bit target, the runtime's conversion of a string
  to a list of characters built the empty list with a boxed \lc{UInt32} zero, which on
  \texttt{wasm32} is a heap object, not a scalar; the patch is three lines and lives in
  \file{engine/wasm/}. It looked like a reference-counting error for an hour.
\end{logbox}

The spike settled the shape of everything after it. The engine would be one program; the
notebook would be thin; the proofs would live next door; and ``verified'' would mean what the
statuses say. The rest of the book is what it took to make the statuses true.

\begin{trybox}[Your first three cells]
\begin{minted}{text}
diff(x^2 * sin(x), x)
rref([1,2,3;4,5,6;7,8,10])
sqrt(50) - sqrt(18)
\end{minted}
  Open the work on each, and click a subterm of an answer. The third cell's answer is $2\sqrt{2}$,
  and Chapter~\ref{ch:cannot} explains why the engine could not write it that way as a term.
\end{trybox}

\chapter{Rewriting with a Trace}\label{ch:rewrite}

Everything the engine does to a term, it does by rewriting: a rule looks at one node, and either
does not apply or replaces the node by something else and says why. The rewriter applies rules
until none applies, records every firing, and hands the record to the notebook. This chapter is
that machinery: the terms, the rules, the strategy, and the trace. It is deliberately small; the
weight of the book is in what the rules \emph{claim} (Part~II) and in why the rewriting
\emph{stops} (Part~III).

\section{Terms}

The term language is deliberately tiny.

\begin{minted}{lean}
inductive Expr where
  | num    : Q → Expr                    -- an exact rational, with an "approximate" flag
  | var    : String → Expr
  | add    : List Expr → Expr            -- n-ary sums and products
  | mul    : List Expr → Expr
  | pow    : Expr → Expr → Expr
  | fn     : String → List Expr → Expr   -- sin, ln, diff, matrix commands, …
  | matrix : List (List Expr) → Expr
\end{minted}

There is no subtraction and no division: $a - b$ is \lc{add [a, mul [-1, b]]} and $a/b$ is
\lc{mul [a, pow b (-1)]}. The printer recovers the familiar forms on the way out, and the
rewriter never has to know about them. Numerals are core Lean's \lc{Rat}, wrapped in a structure
\lc{Q} that adds a flag saying whether the number came from a decimal literal (so that
\texttt{2.5} prints as \texttt{2.5} and not \texttt{5/2}); the flag is presentation and changes no
value, and a two-line theorem in \file{proofs/Proofs/Q.lean} says so.

Sums and products are lists, and their arguments are kept in a canonical order. The function
\lc{canon} sorts the arguments of a sum or a product (products containing a matrix are left
alone: matrix multiplication does not commute). Canonical order runs \emph{silently} inside the
rewriter rather than as a rule: it is applied to every node after its children are normalized,
it is never recorded as a step, and Chapter~\ref{ch:origins} will have to bridge over it when it
traces origins across steps.

\section{Rules}

A rule inspects one node and either declines or returns a result: the new node, an explanation
in Markdown-with-\LaTeX, optionally a nested derivation (row reduction records its row operations
this way), and optionally an error --- a rule may refuse the whole evaluation, as the dimension
check in a matrix product does.

\begin{minted}{lean}
structure RuleResult where
  result : Expr
  explanation : String
  sub : Option Derivation := none
  error : Option String := none
\end{minted}

The rules of the original \lc{simplify} carry one more thing: a proof that they make a measure go
down.

\begin{minted}{lean}
structure Weights where
  w : Expr → Nat
  pos : ∀ e, 1 ≤ w e
  head : ∀ e cs, w (withChildren e cs) = w e   -- a weight looks only at the node's head

structure Rule (W : Weights) where
  name : String
  silent : Bool := false
  apply : Expr → Option RuleResult
  decreasing : ∀ e r, apply e = some r → measure W r.result < measure W e
\end{minted}

\lc{measure W} is the weighted node count. Because a weight depends only on a node's head and is
at least one, two facts are proved once for every weighting: the measure is additive over the
children (replacing children by lighter children makes the node lighter), and it is invariant
under \lc{canon} (a permutation of the arguments weighs the same). Those two facts are what let
the rewriter below be a total, well-founded function rather than a loop with a counter.

\begin{notebox}
  The obligation \lc{decreasing} is discharged by each rule's author, in Lean, at the point of
  writing the rule. For the seven original simplification rules the discharge is a few lines of
  \lc{omega} each under weights $2$ for a numeral, $4$ for a sum or product, $5$ for a power, $8$
  for a variable, a function or a matrix. That was enough for \lc{simplify}. It was not enough for
  derivatives, and Part~III is the story of what replaced it.
\end{notebox}

\section{The strategy}

The rewriter is \emph{innermost}: it normalizes the children of a node first, then applies the
first rule that fires at the node, then normalizes the result again. Innermost is not the only
choice, but it is the one that makes the later termination argument possible
(Chapter~\ref{ch:order}): a rule only ever fires on a node whose children are already normal.

\begin{minted}{lean}
mutual
  def normAt (rules : List (Rule W)) (e : Expr) (path : Path) (acc : Array RawStep) :
      {r : Expr × Array RawStep // measure W r.1 ≤ measure W e} := …
  def normChildren … := …
end

def normalize (rules : List (Rule W)) (e : Expr) : TraceM Expr := do
  let ⟨(out, raw), _⟩ := normAt rules e [] #[]
  let (steps, _) := buildSteps e raw
  modify (· ++ steps)
  pure out
\end{minted}

The return type of \lc{normAt} is a subtype: the function returns its result \emph{together with a
proof} that the result is no heavier than the input. That proof is what the recursion is measured
by --- after a rule fires, the result is strictly lighter, so normalizing it again is a smaller
problem --- and Lean accepts the definition without \lc{partial}, without fuel, and without any
promise the code does not keep.

\section{The trace}

A step records the rule, the explanation, the \emph{path} from the root to the node that changed,
and the whole term before and after.

\begin{minted}{lean}
structure Step where
  rule : String
  explanation : String
  path : Path            -- child indices from the root
  before : Expr
  after : Expr
  sub : Option Derivation := none

structure Derivation where
  input : Expr
  steps : Array Step
  output : Expr
\end{minted}

Paths are how the notebook and the engine refer to the same node. When the engine prints a term as
\LaTeX\ for the notebook it wraps every subterm in an annotation carrying its path,
\verb|\htmlData{path=1.0}{…}|, and KaTeX turns that into a DOM attribute. A click on the rendered
$\cos x$ yields the path \texttt{1.0}; the notebook sends it back with the cell and the line it was
clicked on; the engine walks its own copy of the term to the same node. No text is ever parsed
back. Chapter~\ref{ch:origins} builds on this to answer ``where did this come from?''.

\begin{logbox}[Parity, then a differential test]
  The engine replaced a TypeScript reference implementation. The port was tested against the
  reference at the wire level: 147 cell sources in one session, comparing text, \LaTeX\ with paths,
  value trees, errors and bindings --- zero mismatches, after which the reference was deleted and
  its answers kept as a golden file replayed by \texttt{lake test}. The differential test found
  two things worth keeping: matrix products must not be sorted into canonical order, and the two
  ``parity'' rules $(ab)^n \to a^n b^n$ and $(b^m)^n \to b^{mn}$ with symbolic $m$ resist every
  additive measure. They wait for Chapter~\ref{ch:order}.
\end{logbox}

\section{Exercises}

\begin{exercisebox}[A rule with its obligation]
  Write the rule $x + 0 \to x$ as a \lc{Rule simpW}: the \lc{apply} function on a node
  \lc{.add es}, removing the zero numerals, and the \lc{decreasing} proof. Under which weightings
  does the obligation hold? (It holds for every \lc{Weights}, since a node was removed; say why
  the \lc{head} law is what makes that a one-line proof.)
\end{exercisebox}

\begin{exercisebox}[Silent steps]
  Canonical ordering is not recorded as a step. Give a derivation in which the term after step $n$
  and the term before step $n{+}1$ differ, and describe what the notebook would show if the
  reordering \emph{were} recorded. Chapter~\ref{ch:origins} has to deal with exactly this gap.
\end{exercisebox}

\begin{exercisebox}[Paths]
  The path of the $\cos x$ in $2x\cos x + x^2\sin x$ depends on canonical order. Compute it from
  the tree \lc{add [mul [2, x, cos x], mul [x^2, sin x]]}, then change the order of the two
  summands and compute it again. Why must the engine, not the notebook, be the one to assign
  paths?
\end{exercisebox}

\part{Meaning}
\chapter{Soundness Is a Fold}\label{ch:sound}

A rewrite rule that terminates but changes the value is worse than useless: it is confident and
wrong. This chapter is about the theorems that say a rule \emph{means} something --- first over a
small fragment where everything is decidable, then over the real numbers --- and about the two
rules that turned out not to, whose refutations are the most instructive theorems in the engine.

\section{Soundness for a rule, and for a run}

Fix a relation \lc{rel} between terms that we intend to mean ``these denote the same thing''. A
rule is sound for \lc{rel} if its output is related to its input. That is a statement about one
firing at one node. What the notebook needs is a statement about the \emph{whole run}: the input
of a cell is related to its output. The passage from one to the other is a fold over the
rewriter's recursion, and it needs exactly four properties of the relation.

\begin{minted}{lean}
structure Congruence where
  rel : Expr → Expr → Prop
  refl : ∀ e, rel e e
  trans : ∀ {a b c}, rel a b → rel b c → rel a c
  /-- Rewriting each child soundly rewrites the node soundly. -/
  congr : ∀ e cs, RelList rel (children e) cs → rel e (withChildren e cs)
  /-- Reordering a sum's or product's arguments is sound — what `canon` does silently. -/
  canon : ∀ e, rel e (canon e)

def RuleSoundFor (C : Congruence) (r : Rule W) : Prop :=
  ∀ e res, r.apply e = some res → C.rel e res.result

theorem normalize_sound_for (C : Congruence) (rules : List (Rule W))
    (hs : ∀ r ∈ rules, RuleSoundFor C r) (e : Expr) (s : Array Step) :
    C.rel e ((normalize rules e).run' s)
\end{minted}

The theorem is proved once, for an abstract \lc{Congruence}, by following the rewriter's
recursion: children are rewritten (congruence), the node is put into canonical order
(\lc{canon}), a rule fires (its own soundness), the result is rewritten again (transitivity,
and the recursion). Nothing in the proof knows what \lc{rel} is. That was a deliberate
generalization made in the third week, and it paid three times: the integers, the reals, and
the reals under a derivative each instantiate the same fold.

\section{The integer fragment}

The first semantics is the smallest one that can be \emph{decided}: an evaluator
\lc{eval?}, from an environment and a term to an optional integer, that gives a value to numerals with denominator one,
variables from an environment, sums, products and natural powers, and \lc{none} to everything
else --- functions, fractions, matrices. Two terms are related when the second refines the first:
wherever the first has a value, the second has the same one.

\begin{minted}{lean}
def Refines (e e' : Expr) : Prop := ∀ ρ v, eval? ρ e = some v → eval? ρ e' = some v
theorem simplify_sound (e : Expr) : Refines e (simplify0 e)
\end{minted}

Every one of the seven simplification rules is sound for this relation, and the fold gives
\lc{simplify_sound}. On its own this is a modest claim --- it says nothing about $\sqrt 2$ or
$\sin$ --- but it was the first end-to-end theorem about the engine, checked on the first day the
rules existed, and it fixed the \emph{shape} every later semantics would take.

\section{The reals}

The real semantics lives in \file{proofs/}, because $\mathbb{R}$ is Mathlib's. It is a total
function: powers are \lc{Real.rpow}, \texttt{sqrt} is \lc{Real.sqrt}, \texttt{ln} is
\lc{Real.log}, and where these are undefined in mathematics they take Mathlib's \emph{junk} values
($\log 0 = 0$, $\sqrt{-1} = 0$, $0^{-1} = 0$). Junk values are not a weakness here; they are what
makes every statement below an unconditional equation between total functions, so that the side
conditions a rule really needs have to be \emph{written down} rather than hidden inside a partial
function.

\begin{minted}{lean}
def evalR (ρ : EnvR) : Expr → ℝ
  | .num q => (q.val : ℝ)
  | .var x => ρ x
  | .add es => sumR ρ es
  | .mul es => prodR ρ es
  | .pow b e => (evalR ρ b) ^ (evalR ρ e)
  | .fn f [a] => applyFn f (evalR ρ a)
  | .fn _ _ => 0
  | .matrix _ => 0

def SemEqR (e e' : Expr) : Prop := ∀ ρ, evalR ρ e = evalR ρ e'
\end{minted}

A bridge theorem, \lc{evalR_of_eval?}, says that wherever the integer evaluator is defined the two
agree; so the theorem of the previous section was a restriction of the truth, not a different
claim. The reader will meet this pattern again in Chapter~\ref{ch:deriv}: three semantics, each
shown to be a restriction of the next.

\section{Two rules that are not sound}

Five of the seven rules are sound for \lc{SemEqR} and the fold applies to them. Two are not, and
the interesting part is that this is \emph{proved}.

\begin{thmbox}[\lc{not_collectPowers_soundR}]
  The rule $b^m \cdot b^n \to b^{m+n}$ is not unconditionally sound over $\mathbb{R}$: at $x = 0$
  it rewrites $x \cdot x^{-1}$, which is $0$, to $x^0$, which is $1$.
\end{thmbox}

\begin{minted}{lean}
theorem not_collectPowers_soundR : ¬ RuleSoundR collectPowers := by
  intro hs
  have h := hs (.mul [.var "x", .pow (.var "x") (.num (Q.ofInt (-1)))]) _ rfl (fun _ => 0)
  norm_num [evalR_mul, prodR_cons, prodR_nil, evalR_var, evalR_pow, evalR_num, …] at h
\end{minted}

\begin{thmbox}[\lc{not_functionRules_soundR}]
  The rule $\exp(\ln x) \to x$ is not unconditionally sound: at $x = -1$ the left side is
  $\exp(\ln(-1)) = \exp(\ln 1) = 1$, because \lc{Real.log} is even, while the right side is $-1$.
\end{thmbox}

Both rules are the standard convention of every computer-algebra system, and the engine keeps
them. What changed is the label. The collect-powers rule is \status{conditional}, and the theorem
\lc{collectPowers_soundR_on} says exactly when it is right: whenever the merged base is positive.
The notebook shows that condition next to the step. A reader who has only ever seen $x/x = 1$
printed without comment now sees the $x \neq 0$ that was always there.

\begin{logbox}[On pinning Mathlib]
  The proofs package pins Mathlib at the release tag whose \texttt{lean-toolchain} equals the
  engine's. The first plan had been ``pin to Mathlib master''; on the day, master was on a
  release candidate \emph{ahead} of stable, so following it would have forced the engine off the
  policy of latest stable. Pin to the tag that matches; bump both together.
\end{logbox}

\section{What the fold does not cover}

The fold is stated for rules that are sound \emph{unconditionally}. A conditional rule cannot be
folded, because the hypothesis at one step is about the values at that step, and the fold does
not know them. So the notebook's ``this derivation is sound'' is not one theorem; it is the
conjunction of the statuses of the steps, exactly as displayed. Chapter~\ref{ch:integrate} turns
this into a precise statement --- the integral checker's theorem takes ``the derivation shown was
sound at the point'' as its hypothesis, and that hypothesis is what the statuses say.

\begin{trybox}
\begin{minted}{text}
x * x^(-1)
exp(ln(x))
\end{minted}
  Both simplify. Click each step and read the proof status: \status{conditional}, with the
  hypothesis written out. Then try \texttt{subst(x * x\^{}(-1), x, 0)}: the answer is $1$,
  because the children of a command are normalized first and the conditional step has already
  fired. The label on that step is the only warning, which is the point of the label.
\end{trybox}

\section{Exercises}

\begin{exercisebox}[A third refutation]
  The rule $\sqrt{x^2} \to x$ is not in the engine. Prove in the style of the two refutations
  above that it would not be unconditionally sound over $\mathbb{R}$, and state the side
  condition under which it is.
\end{exercisebox}

\begin{exercisebox}[Congruence laws]
  Show that \lc{SemEqR} satisfies the \lc{canon} law: a permutation of the arguments of a sum
  does not change \lc{sumR}. (The engine's proof is \lc{sumR_perm}, an induction on
  \lc{List.Perm}.) Which of the four laws would fail for a relation that compared terms only at
  the single point $\rho = 0$?
\end{exercisebox}

\chapter{Derivatives and the Binder Problem}\label{ch:deriv}

The derivative rules were the first rules that could not simply be added to the existing story.
Their termination broke the additive measure (Part~III), and their \emph{meaning} broke the
semantics of the previous chapter in a way that is worth understanding, because it is the
reason the engine has three semantics rather than one.

\section{The rules}

$\texttt{diff}(f, x)$ is a function node, \lc{.fn "diff" [f, .var x]}, and the derivative rules
push it inward: the derivative of a sum is the sum of the derivatives, of a product the product
rule, of a power the power rule, of $\sin u$ the chain rule with $\cos u$ outside. Each is one
rewrite, each is recorded, and the trace of $\frac{d}{dx}(x^2 \sin x)$ reads like a textbook
because it \emph{is} the textbook's derivation, one rule per line.

\section{Why \lc{evalR} could not take the case}

The obvious plan was to extend \lc{evalR} with a case for \texttt{diff}. It cannot be done, and
the reason is the second child.

In $\texttt{diff}(f, x)$ the argument $x$ is not a value; it is a \emph{binder}, the name of the
variable we differentiate along. But \lc{Expr} has no binder positions --- every child of every
node is a term. So \lc{.var x} and \lc{.add [.var x]} denote the same real number, while only one
of them can be read as ``the variable we differentiate along''. A semantics that interprets
\texttt{diff} therefore cannot satisfy the \lc{congr} law of Chapter~\ref{ch:sound}: rewriting
the child \lc{.add [.var x]} to \lc{.var x} is sound for every other node and changes the meaning
of this one. And \lc{congr} is the law the fold rests on.

\section{A third semantics}

So the semantics was extended rather than edited. \lc{evalD} agrees with \lc{evalR} everywhere
except that it reads \texttt{diff}, via Mathlib's \lc{deriv}:

\begin{minted}{lean}
def fnD : EnvR → String → List Expr → ℝ
  | ρ, f, [a] => applyFn f (evalD ρ a)
  | ρ, f, [g, v] =>
    if f = "diff" then
      match v with
      | .var x => deriv (fun t => evalD (upd ρ x t) g) (ρ x)
      | _ => 0
    else 0
  | _, _, _ => 0
\end{minted}

The theorem \lc{evalD_eq_evalR} says the two agree on every term without a \texttt{diff} node.
This is the third layer of one pattern: \lc{eval?} $\subset$ \lc{evalR} $\subset$ \lc{evalD},
each proved a restriction of the next. And because the derivative rules never run under the
verified \lc{normalize} of Chapter~\ref{ch:rewrite} --- they run under the tiered rewriter of
Part~III --- they never needed the congruence law in the first place. Each rule gets its own
theorem, at the whole-term level, with \lc{deriv}.

\section{Which rules need a hypothesis}

Here the junk-value discipline of the previous chapter earns its keep. \lc{deriv} is junk-valued
where a function is not differentiable, so the sum, product and chain rules genuinely fail
without a hypothesis, and the hypothesis has to be stated.

\begin{itemize}[nosep]
  \item \status{verified}: \rn{diff.constant}, \rn{diff.variable}, \rn{diff.constant-multiple}.
    Pulling a constant factor out needs no differentiability at all.
    The Mathlib lemma is \lc{deriv_const_mul_field}.
  \item \status{conditional}: \rn{diff.sum} (every summand differentiable at the point),
    \rn{diff.product} (both factors), \rn{diff.power} (a natural exponent and a differentiable
    base; the real-exponent case is open), \rn{diff.chain} (the inner function differentiable;
    proved for $\sin$, $\cos$, $\exp$; $\ln$ and $\tan$ also need a domain condition and are open).
\end{itemize}

And the necessity is proved, as it was for the algebraic rules:

\begin{thmbox}[\lc{not_diff_sum_sound}]
  The sum rule is not unconditional: at $0$, the rule gives $|x| + x$ the derivative $1$, where
  the derivative does not exist.
\end{thmbox}

The engine reports all of this through the capabilities message, so the notebook's panel says
``conditional: needs both factors differentiable at the point'' instead of ``no soundness
theorem yet''. That message was the first place where the statuses of the Preface became
visible to a user.

\begin{logbox}[Not claimed]
  \rn{diff.higher-order} is an abbreviation --- it eliminates the three-argument form
  $\texttt{diff}(f, x, n)$ into $n$ nested derivatives --- so there is nothing to prove.
  \rn{diff.matrix} is proved but vacuous: a matrix has no value in the real semantics, so the
  theorem has no content. The honest label for both is \status{unverified}, and the note next to
  the label says why.
\end{logbox}

\begin{trybox}
\begin{minted}{text}
diff(abs(x) + x, x)
diff(x^x, x)
diff(ln(x), x)
\end{minted}
  The first is the counterexample of \lc{not_diff_sum_sound} in the notebook: the sum rule
  fires, leaves the derivative of $|x|$ standing, and its status says what it assumed. The second uses a rule the theorems do not yet cover
  (a real exponent) and says so.
\end{trybox}

\section{Exercises}

\begin{exercisebox}[The binder problem, concretely]
  Give two terms that \lc{evalR} cannot distinguish but that a semantics for \texttt{diff} must,
  and show which law of \lc{Congruence} they break.
\end{exercisebox}

\begin{exercisebox}[A missing case]
  The chain rule for $\ln u$ needs $u > 0$ at the point (and differentiability). Write the
  theorem statement in the style of \lc{diff_chain_sin_sound}, and say which Mathlib lemma
  about \lc{Real.log} you would need.
\end{exercisebox}

\part{Termination Without a Budget}
\chapter{The Ordering}\label{ch:order}

Chapter~\ref{ch:rewrite} promised that the simplifier would never need a step budget, and kept the
promise with an additive measure. This chapter is about the day that measure met the rest of the
notebook --- derivatives that duplicate subterms, commands that grow terms, matrix rules that
copy a scalar into every entry --- and about the ordering that replaced it. It is the chapter
where the project's order theory is used rather than displayed, and the one whose constants were
hardest won.

\section{Why the additive measure fails}

The measure of Chapter~\ref{ch:rewrite} is a weighted node count, and it is additive: the measure
of a node is a head weight plus the measures of its children. Additivity is exactly what the fold
over the rewriter needs, and it is exactly what a duplicating rule breaks. The product rule sends
$\frac{d}{dx}(f \cdot g)$ to $f' g + f g'$: two copies of $f$ and two of $g$ where there was one
of each. No head weight for \texttt{diff} can pay for that, because the payment must be
\emph{proportional to the size of $f$}, and an additive weight is a constant.

The classic answer is a recursive path ordering, where a symbol higher in a precedence may be
replaced by any term built from lower ones. It fits the derivative rules perfectly --- they are the
textbook example --- and it does not fit the rest of the simplifier. Two rules already in the
engine want opposite precedences:

\[
  (ab)^n \to a^n b^n \quad\text{needs}\quad \mathtt{pow} > \mathtt{mul},
  \qquad
  x \cdot x \to x^2 \quad\text{needs}\quad \mathtt{mul} > \mathtt{pow}.
\]

Neither could be moved. Distributing integer powers over products is what every computer-algebra
system does under ``simplify'', and $x \cdot x \to x^2$ is the reason collect-powers exists. A
global RPO was also not in Mathlib, and the estimate for building one was three times the work of
the alternative --- while still leaving these two rules unordered.

\begin{logbox}[The decision]
  Option B, recorded on 13 September with Alex: tiered measures with an innermost-strategy
  lemma. Not a global RPO; not ``fuel as a proven bound''. The two parity rules stay in
  \texttt{simplify} and get a bespoke ordering. Expansion --- whose rules genuinely grow terms ---
  stays on fuel for now and gets its own chapter (\ref{ch:cannot}).
\end{logbox}

\section{Six tiers}

The ordering is lexicographic on a tuple of six natural numbers. Each tier exists because a
family of rules could not be ordered by the tiers below it.

\begin{minted}{lean}
abbrev Mu := Nat × Nat × Nat × Nat × Nat × Nat
def μ (e : Expr) : Mu := (cmdCount e, d3Count e, litCount e, M e, size e, numCount e)
\end{minted}

\begin{enumerate}[nosep]
  \item \textbf{Commands.} \texttt{expand}, \texttt{rref}, \texttt{N}, \texttt{subst},
    \texttt{integrate}: a command node is consumed by its rule and produces a term with no command
    in it. Whatever the command does inside, this tier drops.
  \item \textbf{Malformed derivatives.} The three-argument $\texttt{diff}(f, x, n)$ is eliminated
    into $n$ nested two-argument ones, which \emph{raises} the weight $M$; counting three-argument
    nodes first makes the elimination a decrease.
  \item \textbf{Nodes above matrix literals.} Scalar multiplication $s \cdot A$ copies $s$ into
    every entry, so counting literals alone fails and counting entries fails. Counting the nodes
    \emph{on a path from the root to a literal} makes every matrix rule a strict decrease, and
    keeps matrices out of $M$ entirely.
  \item \textbf{The weight $M$.} The tier that orders derivatives against the simplifier. Its
    definition is the next section.
  \item \textbf{Size.} For rules that keep $M$ and drop a node.
  \item \textbf{Integer magnitudes.} Added later, for the radical rules (Chapter~\ref{ch:cannot}).
\end{enumerate}

Lexicographic products of well-founded orders are well-founded, and the whole proof of that is one
line of \lc{Prod.lex}. The lemma \lc{muLt_of} builds a decrease tier by tier: each tier is
$\le$ under the assumption that the earlier tiers are equal, and the first strict one wins.

\section{The weight $M$}

\begin{minted}{lean}
def M : Expr → Nat
  | .num q => if q.isOne then 1 else if q.isInt then 2 else 4
  | .var _ => 10
  | .add es => ML es + (if es.isEmpty then 3 else 0)
  | .mul es => 2 + ML es + 2 * es.length
  | .pow b e => (M b + 1) * M e - 1
  | .fn "diff" [g, t] => 3 ^ (M g + 3) + M t
  | .fn _ es => 5 * ML es + 10
  | .matrix rows => 10 + MR rows
\end{minted}

Every constant here is forced by a named rule, and the log records which.

\begin{itemize}[nosep]
  \item \textbf{A numeric exponent multiplies its base.} $M(b^e) = (M b + 1) \cdot M e - 1$. This is
    the shape that makes $(b^m)^n \to b^{mn}$ decrease with symbolic $m$ while $x \cdot x \to x^2$
    also decreases. Every polynomial interpretation tried before it failed on one or the other.
  \item \textbf{A per-factor charge on products.} $M(\prod e_i) = 2 + \sum (M e_i + 2)$. The
    $+2$ per factor is what $x \cdot x \to x^2$ spends: two factors at 10 each plus their charges
    outweigh $(10+1) \cdot 2 - 1 = 21$.
  \item \textbf{One weighs 1.} So that $x \cdot x^a \to x^{a+1}$ decreases: the exponent $a+1$
    costs one more than $a$, and the product's charge pays for it.
  \item \textbf{Non-integer numerals weigh 4.} So that $\sqrt 2 + \sqrt 2 \to 2 \sqrt 2$
    decreases: the term $2^{1/2}$ must weigh at least $9$, and a fraction weighing $2$ would make
    it $8$.
  \item \textbf{A derivative is exponential in its body.} $M(\texttt{diff}(f, x)) = 3^{M f + 3} +
    M x$. The product rule replaces one derivative of $fg$ by two derivatives of the parts plus
    two copies of the parts, and $3^{M f + M g + \ldots}$ dominates $3^{M f + 3} + 3^{M g + 3} +
    2(M f + M g) + \text{charges}$ with room to spare. Three is the smallest base that works; two
    fails on the chain rule.
\end{itemize}

The tuning was the real work of the milestone. A rule that fails to decrease does not fail
loudly; it fails as an \lc{omega} goal that is false, and finding \emph{which} constant to move,
without breaking the thirty lemmas already proved, took most of two days.

\section{The obligation must be conditional}

The product rule duplicates its body. If the body could still contain a command --- an
\texttt{expand} that will grow when it fires --- then no ordering of this shape survives, because
the duplicated command grows twice. So the termination obligation for a rule is not ``its output
is smaller than its input'' but ``its output is smaller than its input \emph{when the input's
children are already normal}'':

\begin{minted}{lean}
structure Ordered (rules : List PlainRule) : Prop where
  decreasing : ∀ r ∈ rules, ∀ e res, ChildrenNormal rules e →
    r.apply e = some res → res.error = none → MuLt (μ res.result) (μ e)
\end{minted}

The rewriter's innermost strategy makes the hypothesis true at every firing --- children are
normalized before the parent is tried --- and the proof carries \lc{Normal} through the recursion
in the \emph{result type} of \lc{normAtT}, together with the invariant $\mu(\text{out}) \le
\mu(\text{in})$ that a parent needs to know about its rewritten children. This is the
``innermost-strategy lemma'' of the decision, and it is where the order theory and the rewriting
strategy meet: the ordering is not on terms but on terms whose children are in normal form.

\section{Honest boundaries}

Three families of rules run subsystems the proof does not see inside: commands (which call the
expander, elimination, floating point), and the matrix rules (unverified arithmetic until
Chapter~\ref{ch:linalg}). For each, the proof needs one property of the output and nothing else,
and that property is \emph{checked at run time}:

\begin{minted}{lean}
def checked (r : RuleResult) : RuleResult :=
  if r.error.isSome then r
  else if cmdCount r.result = 0 then r
  else refuse "a command produced a term that still contains a command"
\end{minted}

A failure surfaces as an error in the cell, never as non-termination and never as a silent wrong
answer. This is the shape --- decide at run time exactly what the proof needs, and refuse
otherwise --- that the rest of the book keeps returning to. It is a boundary, not a budget: no
number is chosen, and the check is a property of the output, not a count of the steps.

Behaviour changed in the process, all of it deliberately and all of it logged: \texttt{det} and
$A^k$ on literals are one step each; $N(1/3)$ evaluates where the reference left it; $\sin(A)$ for a
matrix $A$ is an error instead of a junk normal form; $\texttt{diff}(f, 2)$ is an error.

\begin{trybox}
\begin{minted}{text}
diff(x^2 * sin(x) * exp(x), x)
(a*b)^3 * (a*b)^2
\end{minted}
  The first is a long derivation with a product rule inside a product rule; there is no budget
  behind it. The second fires both parity rules in sequence. Open the capabilities panel: the
  termination line names \lc{pipelineOrdered}.
\end{trybox}

\section{Exercises}

\begin{exercisebox}[Forcing a constant]
  Compute $M$ on both sides of $x \cdot x \to x^2$ and of $x \cdot x^a \to x^{a+1}$ with $M(1)$
  set to $2$ instead of $1$. Which rule stops decreasing?
\end{exercisebox}

\begin{exercisebox}[Why not base 2]
  Write the chain rule $\frac{d}{dx} \sin u \to \cos u \cdot \frac{d}{dx} u$ and compute $M$ on
  both sides with $M(\texttt{diff}(f,x)) = 2^{M f + 3} + M x$. Find a $u$ for which it fails to
  decrease, then show that base 3 succeeds for every $u$.
\end{exercisebox}

\begin{exercisebox}[The innermost lemma]
  Give a rule set and a term for which the outermost strategy fails to terminate but the
  innermost one terminates under \lc{Ordered}. (A command inside the body of a product is the
  engine's own example.)
\end{exercisebox}

\chapter{What the Ordering Cannot Afford}\label{ch:cannot}

An ordering is a budget of a different kind: it says which rewrites are affordable, and a rule the
ordering cannot afford is not a rule the engine can have. This chapter is about three rules that
were wanted and could not be afforded as rewrites --- expansion, radical simplification, and one
trigonometric identity --- and about the three different things that were done instead. Each of
them has a lesson that the ordering chapter alone would not teach.

\section{Expansion is a function, not a rule set}

Distribution $a(b + c) \to ab + ac$ grows a term, and it grows it in a way no interpretation
compatible with $x \cdot x \to x^2$ can pay for. The proof that it cannot is short and was the
first thing tried: an interpretation under which distribution decreases makes the product $ab$
cheaper than the sum $a + b$ scaled by the multiplier, and that same interpretation makes
$x \cdot x \to x^2$ increase. It is the wall of Chapter~\ref{ch:order} again, seen from the other
side.

So expansion is not a rewrite system. \lc{Expand.dist} is one total, structurally recursive
function: children first, then a product with a sum among its factors is multiplied out, collecting
like monomials as the product is built factor by factor. The pipeline of Chapter~\ref{ch:order}
then collects whatever is left.

\begin{minted}{lean}
def distMul (fs : List Expr) : Expr :=
  if fs.any isAdd
    then ofPoly ((fs.map toPoly).foldl polyMul [⟨Q.one, []⟩])
    else .mul fs
\end{minted}

The collection during distribution is not a nicety. The first version distributed and left the
collection to the pipeline's collect-like-terms rule, which took the 4096 monomials of
$(x+y)^{12}$ one pair at a time at the root, with a full canonical sort after each pair: minutes,
and a recorded derivation of 49\,143 steps. Collecting as a polynomial --- sorted keys, coefficients
folded --- makes it instant and cuts the recorded expansion to two steps.

The soundness theorem is unconditional, because distribution holds in every commutative ring:

\begin{minted}{lean}
theorem dist_sound (ρ : EnvR) : ∀ e : Expr, evalR ρ (Expand.dist e) = evalR ρ e
\end{minted}

It is proved once for an abstract semantics --- a structure \lc{Sem} with an evaluator and its
four equations for sums, products, numerals and natural powers --- and instantiated twice, for
\lc{evalR} and for \lc{evalD}. The second instance matters in Chapter~\ref{ch:integrate}.

\begin{logbox}[Two shapes that read the same]
  A step-recording algorithm and a rewrite system look identical in the notebook: a list of
  steps, each with a rule name and a before and after. That is the argument for writing a
  function whenever a joint ordering would be forced. The reader loses nothing; the proof
  becomes an induction on the function instead of a fold over rules; and the fuel rewriter,
  which expansion had been the last user of, is deleted. There is no step budget anywhere in the
  engine except one, and that one is Chapter~\ref{ch:worlds}'s.
\end{logbox}

\section{Radicals: what the ordering allows}

Alex asked for $\sqrt 8 \to 2\sqrt 2$, and asked for it by retuning $M$. The analysis showed that
retuning cannot do it, and the log records both the request and the outcome.

As a term, $2\sqrt 2$ is $2 \cdot 2^{1/2}$, whose weight is $19$; $8^{1/2}$ weighs $11$. The result
contains the input's radical \emph{and more}. Only a weight that depends on the \emph{value} of a
numeral could pay for that --- making $8$ heavier than $2$ by enough --- and numeral weights must
stay bounded, or fold-constants stops decreasing on arbitrary sums ($1 + 7 \to 8$ would gain).

What the ordering does allow is the single-power form: $8^{1/2} \to 2^{3/2}$, at equal $M$ and
equal size. To order it, a sixth tier was added: the magnitudes of the integer numerals, counted
\emph{head-only}. That last word cost a day. A tier that reads a child's value fails the
children-monotonicity lemma the rewriter rests on --- a numeral child of equal weight could be
replaced by another --- and putting the value on the numeral node itself sidesteps it. The other
constraint is that $4^{2/3} \to 2^{4/3}$ ties every bit-length or value weight on the exponent
side, which is why the tier weighs integers only and non-integer exponents nothing.

Two more rules then do what Mathematica does where it matters. Radicals with the same square-free
part collect in a sum ($\sqrt{50} - \sqrt{18} \to 2\sqrt 2$; a sum of two radicals outweighs one
product, so this decreases $M$), and same-index radicals multiply under one root ($\sqrt{12} \cdot
\sqrt 3 \to \sqrt{36} \to 6$). Each rule guards itself with the decidable decrease its proof needs,
so the ordering lemma is a split on the guard; the guards never fail in practice and are honest
boundaries, not fuel. All three are proved sound over $\mathbb{R}$ unconditionally, their bases
being positive integers.

The printer then shows the single-power form the textbook way: $2^{3/2}$ as $2\sqrt 2$,
$12^{1/2}$ as $2\sqrt 3$, $5\sqrt{12}$ as $10\sqrt 3$. Text output stays faithful to the term. The
notebook shows what a reader expects; the engine holds what the ordering can carry; and the
distance between them is one function in the printer.

\section{Sine over cosine is tangent}

The third rule was found by the integration checker (Chapter~\ref{ch:integrate}) refusing a
correct answer: $\int \tan x$ was rejected because $\frac{d}{dx}(-\ln \cos x)$ normalized to
$\sin x \cdot (\cos x)^{-1}$ and the engine did not identify that with $\tan x$. The rule is a new
case of \rn{simp.function}, it \emph{does} decrease --- one function node replaces a product of two
--- and it had to be proved in all three layers: the additive measure of Chapter~\ref{ch:rewrite},
the ordering of Chapter~\ref{ch:order}, and soundness over $\mathbb{R}$, which is unconditional
because Mathlib's $\tan$ is $\sin/\cos$ everywhere, junk values included.

The lesson is about the cost of a rule. Three layers of proof for one identity is the price of
every rule in the engine, and it is why rules are added when a checker refuses something and not
before.

\begin{trybox}
\begin{minted}{text}
expand((x + y)^12)
sqrt(50) - sqrt(18)
sqrt(12) * sqrt(3)
\end{minted}
  The first has two expansion steps; everything after them is the pipeline collecting. The second and third run the radical rules; click a step and
  see the theorem name.
\end{trybox}

\section{Exercises}

\begin{exercisebox}[The wall, from the other side]
  Suppose an interpretation $I$ makes $a(b+c) \to ab + ac$ strictly decrease for all $a, b, c$.
  Show that $I(x \cdot x) < I(x^2)$ cannot also hold with $I$ of the form in
  Chapter~\ref{ch:order}.
\end{exercisebox}

\begin{exercisebox}[Why head-only]
  Give a rule that is a decrease under a value-of-children tier and a term on which the
  children-monotonicity lemma ($\mu(\text{child}') \le \mu(\text{child}) \Rightarrow
  \mu(\text{parent}') \le \mu(\text{parent})$) fails for that tier.
\end{exercisebox}

\part{Showing the Work}
\chapter{Origins}\label{ch:origins}

``Show work'' has two halves. The first is a derivation the reader can follow line by line, and
Part~I built it. The second is the answer to a question about a \emph{piece} of the result:
where did this $2$ come from? which steps produced $\cos x$? The first version of the engine
answered by comparing paths --- if the selected subterm's path was a prefix of a step's path, the
step was involved --- and it was wrong often enough that a milestone was spent replacing it.

\section{Origin tracking}

The idea is van Deursen, Klint and Tip's (1993): trace a position \emph{backwards} through the
derivation, one step at a time, and let each step say what it did to it.

\begin{itemize}[nosep]
  \item \textbf{Outside the redex}, a position is its own origin. The step did not touch it.
  \item \textbf{Inside the contractum}, the origins of a position are the positions of
    \emph{equal subterms of the redex}: the rule \emph{copied} the node there.
  \item \textbf{With no equal subterm}, the step \emph{created} the node. Its children are traced
    on; the node itself has no earlier history.
  \item \textbf{A position that contains the redex} has the step below it: the relation is
    \emph{contains}.
\end{itemize}

\begin{minted}{lean}
inductive Relation where
  | created   -- the node at the position was built by the rule
  | copied    -- the rule moved or duplicated it into the contractum
  | contains  -- the rule fired strictly below the node
\end{minted}

The result is one relation per step, in derivation order. It is sent to the notebook as a trace,
and \texttt{explain} accepts which term the path is into --- the input, the output, or the term
after step $n$ --- so every subterm of every line is selectable.

\section{The theorem that lets a step be skipped}

The first bullet is the one that needs a proof. Skipping a step ``because it did not touch the
position'' is only sound if a rewrite at $p$ really leaves every position disjoint from $p$
unchanged, and that is a statement about \lc{replaceAt} and \lc{at?}, the two path primitives:

\begin{minted}{lean}
def Disjoint (p q : Path) : Prop := ¬ p.IsPrefix q ∧ ¬ q.IsPrefix p

theorem at?_replaceAt_disjoint (new : Expr) : ∀ (e : Expr) (p q : Path), Disjoint p q →
    (replaceAt e p new).at? q = e.at? q
\end{minted}

It is a small theorem, but it is the one that makes the tracer's cheapest case correct without
looking at the term, and it was written before the tracer.

\section{Three findings}

\begin{enumerate}[nosep]
  \item \textbf{A trace must stop at a created node but continue into its parts.} The product
    rule creates the sum $f'g + fg'$; the node is new, but $f$ and $g$ inside it are copies. If
    the trace stopped at the creation, the history of $2x \cos x$ would be ``the product rule''
    and nothing else. The tracer stops the node and traces its children.
  \item \textbf{Exact equality breaks at the first silent sort.} The rewriter puts arguments in
    canonical order between recorded steps, so $x^{3 + -1}$ in one step's \emph{after} is
    $x^{-1 + 3}$ in the next step's \emph{before}. Equality is therefore taken up to the argument
    order that \lc{canon} may change (\lc{canonDeep}), and the same comparison bridges the silent
    reorderings between recorded steps. The first test caught this.
  \item \textbf{What remains over-approximate.} Equal subterms at several positions --- the
    paper's \emph{secondary origins} --- are reported as sets, never dropped. A reader who selects
    one of two identical $x$'s gets both histories, which is the honest answer.
\end{enumerate}

\begin{trybox}
\begin{minted}{text}
diff(x^2 * cos(x), x)
\end{minted}
  Click the $2x$ in the result. The panel shows the steps that produced it, each labelled
  created, copied or contains. Now click a step, then a subterm of \emph{that} line.
\end{trybox}

\section{Exercises}

\begin{exercisebox}[Disjointness]
  Prove that if \lc{Disjoint p q} then \lc{Disjoint (i :: p) (i :: q)}, and explain which case
  of \lc{at?_replaceAt_disjoint} needs it.
\end{exercisebox}

\begin{exercisebox}[Secondary origins]
  Find a derivation in which a selected subterm has two origins in the input, and say what a
  single-origin tracer would have to invent to report only one.
\end{exercisebox}

\chapter{The Notebook}\label{ch:notebook}

The notebook evaluates nothing. That is its whole design, and it was decided before a line of it
was written: the frontend owns no printer, no simplifier, no numeric evaluator; it sends source to
an engine and draws what comes back. This chapter is about what that leaves for the notebook to
do, which turns out to be more than expected, and about the three ways the engine reaches it.

\section{One engine, three hosts}

The engine is a Lean program with a JSON-RPC surface (Chapter~\ref{ch:what}). It runs in three
places from one source:

\begin{itemize}[nosep]
  \item \textbf{In the browser}, compiled to WebAssembly and loaded in a web worker. The runtime
    and \lc{Init} are built from source for \texttt{wasm32}, because Lean stopped shipping a
    prebuilt one after v4.15 and the project wants current Lean. The module is 1.8~MB and
    answers its first request 96~ms after the worker starts. No threads, so no cross-origin
    isolation headers: a self-hosted notebook is a static directory.
  \item \textbf{Natively}, behind an HTTP host, for a shared server.
  \item \textbf{On the command line}, over stdio, which is how the tests drive it.
\end{itemize}

The notebook does not know which. It holds an \lc{EngineClient} built from a \lc{Transport},
and a transport moves strings.

\section{What the notebook draws}

Everything the notebook shows is either a string from the engine or a picture described by it.
The engine renders terms to text and to LaTeX, and with paths requested wraps every subterm in an
\texttt{\textbackslash htmlData} annotation; the notebook typesets the LaTeX and attaches a click
handler that reads the path back. That is the whole of selection.

The derivation panel is a list of steps. A step click shows the trail up to it
(Chapter~\ref{ch:origins}); a nested derivation --- the simplification inside a derivative step,
or the differentiation inside an integration check --- opens as a sub-panel with the same
machinery, and the status shown for a step is the weakest status of anything beneath it.

Pictures follow protocol rule 5: a new capability is an optional field, never a change. A plot is
a list of samples with \texttt{null} where the function has no finite value; the notebook draws the
curve, breaks it at the gaps, and trims the tails of the range so an asymptote does not flatten
the rest. A Hasse diagram is a list of elements with heights and a list of covers; the notebook
lays them out in layers. In both cases the engine simplified the formula or decided the order,
recorded a derivation, and described the picture; nothing in the drawing evaluates.

\section{Files, and the session behind them}

A notebook file (\texttt{.chalk}) is JSON: the cells' sources, their saved outputs and derivations, and the studio's
scenes. Opening one shows the saved outputs at once and then re-runs every cell in order, so that
the engine's session --- and with it \texttt{explain}, which needs the derivation to be in the
engine's memory --- matches what is on screen. The notebook also autosaves to the browser after
every run and comes back on reload the same way.

\begin{logbox}[The stale module]
  A browser served last week's WebAssembly for a whole afternoon while the tests passed. Every
  asset URL now carries a build stamp, and the notebook logs the engine's version next to the
  worker's start-up time. Not a finding about mathematics; a finding about trust.
\end{logbox}

\section{Input}

Two small conveniences turned out to matter for the worlds of Chapter~\ref{ch:worlds}. Typing a
backslash opens a completion list of every symbol abbreviation, filtered as you type; at
\texttt{\textbackslash lam} it shows the one entry with the symbol it produces, $\lambda$, on the
right, and Tab or a following space inserts it. The same table covers \texttt{\textbackslash pi},
\texttt{\textbackslash e} (Euler's number, which the engine reads as $\exp 1$ so every rule about
$\exp$ applies to it), \texttt{\textbackslash phi} and the Greek alphabet. This is the input method of the Lean editor
extension, borrowed because the reader of this book already knows it. And a cell's kind --- an
ordinary term, a $\lambda$-term, an order-world command --- is decided by the engine and reported
in the reply; the notebook only shows the badge.

\section{Output references and forms}

Two conveniences borrowed from Mathematica. \texttt{\%} is the previous output, \texttt{\%\%} the
one before, \texttt{\%n} the output of evaluation $n$; the engine numbers every evaluation and
reports the number, the notebook shows it as \texttt{In[n]} and \texttt{Out[n]}, and the session
substitutes the referenced value before parsing goes any further, so the input interpretation shows
what \texttt{\%} stood for. And every output carries a small form chip: a matrix can be shown with
square or round brackets, as a bare grid, as a table, or in the engine's input syntax. The engine
prints once; the form is a typesetting choice the notebook keeps in the file.

\section{Exercises}

\begin{exercisebox}[A fourth host]
  The transport for a WebSocket is under thirty lines. Write it, and say what the engine has to
  change. (Nothing.)
\end{exercisebox}

\begin{exercisebox}[Rule 5 in practice]
  Add a field to \lc{EvaluateResult} that an old notebook would ignore and a new one would draw,
  and explain why the old notebook keeps working against the new engine and the new notebook
  against the old one.
\end{exercisebox}

\part{Three More Worlds}
\chapter{Linear Algebra over $\mathbb{Q}$}\label{ch:linalg}

Row reduction was in the engine from the first week, as a step-recording algorithm driven by the
simplifier, and for months its steps were labelled \status{unverified}. This chapter is about the
two theorems that changed the label: elimination preserves the solution set, and the result is in
reduced row echelon form. Both live in the engine itself, with no Mathlib, and the second cost a
lemma that reads better than the algorithm it is about.

\section{A matrix is a system}

The first decision was what a matrix \emph{means}. Reading it as a linear map would have needed
matrix multiplication and its algebra; reading it as a homogeneous system needs a dot product.
One equation $r \cdot x = 0$ per row, and the augmented matrix of $Ax = b$ is the same system at
the point $(x, -1)$, so both readings are one predicate:

\begin{minted}{lean}
def Sol (m : Mat) (x : List Rat) : Prop := ∀ r ∈ m, dot r x = 0
\end{minted}

Elimination is a list of the three elementary row operations, and the theorem is one lemma per
operation folded over the list:

\begin{minted}{lean}
inductive Op where
  | swap (i j : Nat)              -- exchange rows i and j
  | scale (i : Nat) (c : Rat)     -- Rᵢ ← c · Rᵢ; with c = 0 it is the identity
  | addMul (i j : Nat) (c : Rat)  -- Rᵢ ← Rᵢ + c · Rⱼ; with i = j it is the identity

theorem sol_rref (m : Mat) (x : List Rat) : Sol (rref m) x ↔ Sol m x
\end{minted}

\section{Three ways to have no side condition}

The lemmas have no hypotheses, and each absence was arranged.

\begin{enumerate}[nosep]
  \item \textbf{Degenerate parameters are the identity.} Scaling a row by $0$ would lose an
    equation; adding a row to itself with coefficient $-1$ would too. Instead of proving that
    the algorithm avoids these cases, the operations \emph{define} them as doing nothing. Then
    every operation is invertible for free, and the algorithm needs no proof that it stays
    away from the edge.
  \item \textbf{Pad, do not truncate.} Adding a shorter row to a longer one pads the shorter with
    zeros (\lc{rowAdd}). This removes the rectangularity hypothesis from every lemma, because
    nothing ever depends on rows having equal length.
  \item \textbf{The field algebra is \lc{grind}'s.} Mathlib's \lc{ring} is not available in the
    Init-only engine, and the identities over \lc{Rat} that each lemma needs are discharged by
    \lc{grind}, whose ring and field instances ship with core. So the whole proof lives in
    \file{engine/}, next to the code it is about.
\end{enumerate}

The command's numeral path replays the emitted operations into the derivation steps, so the
notebook shows \rn{la.row-swap}, \rn{la.row-scale} and \rn{la.row-add} with the theorem names.
Symbolic entries keep the old simplifier-driven algorithm under a \texttt{.symbolic} suffix and
stay \status{unverified}, because a symbolic pivot is only \emph{assumed} nonzero.

\section{Echelon form: proved, not checked}

The first version decided at run time that the output was in echelon form and refused otherwise
--- the boundary of Chapter~\ref{ch:order}, applied once more. It was on the open list because a
boundary is a promise to come back. The proof carries an invariant column by column:

\begin{minted}{lean}
/-- Rows `< p` are pivot rows with pivot columns `pivs`; rows `≥ p` are zero left of `col`. -/
structure Inv (m : Mat) (p col : Nat) (pivs : List Nat) : Prop where
  len    : pivs.length = p
  lt     : ∀ c ∈ pivs, c < col
  sorted : pivs.Pairwise (· < ·)
  pivot  : ∀ i, i < p → entry m i (pivs.getD i 0) = 1
  lead   : ∀ i, i < p → ∀ j, j < pivs.getD i 0 → entry m i j = 0
  alone  : ∀ i, i < p → ∀ k, k ≠ i → entry m k (pivs.getD i 0) = 0
  below  : ∀ i, p ≤ i → ∀ j, j < col → entry m i j = 0
  p_le   : p ≤ m.length

def IsRref (m : Mat) (w : Nat) : Prop := ∃ p pivs, Inv m p w pivs
theorem rref_isRref (m : Mat) : IsRref (rref m) (ncols m)
\end{minted}

A column step either finds no pivot --- the invariant moves one column right --- or swaps, scales
and clears, giving one more pivot row. The lemma that made the proof entry-wise rather than
operational is \lc{entry_clears}: the clearing step reads all its factors from the matrix
\emph{before} any clearing, so its sequential application equals the simultaneous one, and each
row changes by its own multiple of the pivot row, which is itself untouched. Without that
observation the proof would have had to follow the loop; with it, every clause of \lc{Inv} is
checked on one entry at a time.

The statement is in the width of the \emph{input}. That sidesteps proving that row operations
preserve rectangularity; for the rectangular matrices the parser admits, the two widths coincide.

\begin{logbox}[Still open]
  Determinants and matrix products over $\mathbb{Q}$, and a value for matrices in the real
  semantics --- which is what would give \rn{diff.matrix} content. None of it is hard; all of it
  is listed, with the honest label, in Appendix~\ref{app:rules}.
\end{logbox}

\begin{trybox}
\begin{minted}{text}
rref([1,2,3;4,5,6;7,8,10])
rref([a,1;1,a])
\end{minted}
  The first is verified all the way down. The second is symbolic: its steps carry the
  \texttt{.symbolic} suffix and say which pivot was assumed nonzero.
\end{trybox}

\section{Exercises}

\begin{exercisebox}[Identity by definition]
  Show that \lc{sol_scale} would need the hypothesis $c \neq 0$ if scaling by $0$ zeroed the
  row, and that the algorithm as written never scales by $0$ anyway. Which of the two is the
  cheaper thing to know?
\end{exercisebox}

\begin{exercisebox}[The width]
  Give a non-rectangular \lc{Mat} for which \lc{IsRref (rref m) (ncols m)} holds but the
  intuitive ``reduced echelon form'' does not, and explain why the parser never produces it.
\end{exercisebox}

\chapter{Integration as a Checker}\label{ch:integrate}

Nobody proves an integrator correct. The engine does not try. What it does instead is the oldest
trick in the subject --- differentiate the answer and see --- made into a theorem, and the
interesting part is what had to change in the pipeline for the theorem to be stateable at all.

\section{A guess, then a check}

The finder (\file{Antiderivative.lean}) is an ordinary table-driven antidifferentiator: sums,
constant factors, powers, $a^u$, the elementary functions, each under a linear substitution;
later, $u$-substitution and integration by parts with the LIATE heuristic. Nothing about it is
proved, and its steps are labelled with a fourth status, \status{checked}: a guess whose result a
later step verifies.

The command then differentiates the candidate with the pipeline, and accepts only if the normal
form of the derivative is the normal form of the integrand, exactly. The claim is a theorem about
the command's output, and it assumes nothing about the finder:

\begin{minted}{lean}
theorem cmdIntegrate_spec (norm : Norm) {f : Expr} {x : String} {res : RuleResult}
    (h : (cmdIntegrate norm).apply (.fn "integrate" [f, .var x]) = some res)
    (hok : res.error = none) :
    ∃ g sub g' s₁ s₂, norm (D res.result x) = .ok (g, sub) ∧
      norm (Expand.dist g) = .ok (g', s₁) ∧ norm (Expand.dist f) = .ok (g', s₂)
\end{minted}

And in \file{proofs/}, with the real semantics, that becomes a statement about \lc{deriv}:

\begin{minted}{lean}
theorem integrate_deriv (norm : Norm) … (ρ : EnvR)
    (hdiff  : … evalD ρ (D res.result x) = evalD ρ g)
    (hleft  : … evalD ρ (Expand.dist g) = evalD ρ g')
    (hright : … evalD ρ (Expand.dist f) = evalD ρ f') :
    deriv (fx ρ x res.result) (ρ x) = evalD ρ f
\end{minted}

The three hypotheses say that the normalizations the check performed were sound at the point
$\rho$. That is exactly what the statuses of the nested derivation say, rule by rule
(Chapters~\ref{ch:sound} and~\ref{ch:deriv}). So the notebook's display of an integration ---
finder steps marked \status{checked}, an \rn{int.check} step with the differentiation nested
under it, and the command inheriting the weakest status below it --- is a faithful rendering of
this theorem's hypotheses.

\section{The knot}

The pipeline could not check with itself. A rule set cannot contain a rule that normalizes with
that set --- the definition would be circular, and Lean says so. The way out was to make the
pipeline \emph{generic in the checker's normalizer}:

\begin{minted}{lean}
def pipelineRulesWith (norm : Norm) : List PlainRule := …
theorem pipelineOrderedWith (norm : Norm) : Ordered (pipelineRulesWith norm)
\end{minted}

The termination theorem of Chapter~\ref{ch:order} then holds for \emph{every} \lc{norm}, and the
knot closes after the proof: the checker is the pipeline with nested \texttt{integrate} refused,
the notebook's pipeline is the pipeline with that checker, and neither has a step budget.

\begin{minted}{lean}
def pipelineRules : List PlainRule := pipelineRulesWith checkNorm
theorem pipelineOrdered : Ordered pipelineRules := pipelineOrderedWith checkNorm
\end{minted}

\section{Exact means refusable}

The check is exact, so a correct guess can be refused. $\int \tan x$ was: the finder produced
$-\ln \cos x$, its derivative normalized to $\sin x \cdot (\cos x)^{-1}$, and the engine did not
identify that with $\tan x$. That is the right failure mode --- never a wrong answer --- and the fix
was a rule (Chapter~\ref{ch:cannot}), proved in three layers, not a weaker check.

Integration by parts found a subtler one. Every candidate was correct and every one was refused,
because the pipeline never distributes a numeral over a sum: the ordering forbids it, distribution
is \texttt{expand}'s job, and so $-(a + b) + b$ is a normal form. The fix keeps the check sound
rather than weakening it: both sides are expanded with \lc{Expand.dist} --- proved sound for the
derivative semantics too, \lc{dist_soundD}, the second instance of the generic proof --- and
simplified before comparison. That is the \rn{int.compare} step, and it is why the spec above
compares \lc{norm (Expand.dist g)} with \lc{norm (Expand.dist f)} rather than \lc{g} with
\lc{f}.

A third finding came free. $\int x^a$ is accepted through \rn{simp.collect-powers} cancelling
$(a+1)/(a+1)$, so its check inherits that rule's side condition, $a \neq -1$, and the statuses say
so without anyone writing it down. The theorem's hypotheses are the notebook's display; the
display knew about $a = -1$ before the author did.

\begin{logbox}[Still refused]
  $e^x \sin x$ needs the ``solve for the integral'' trick; $\sin^2 x$ needs a trigonometric
  identity the simplifier lacks; $1/(x^2 + 1)$ needs an arctangent the engine does not have. Each
  is refused with the candidate it tried, and none is a wrong answer.
\end{logbox}

\begin{trybox}
\begin{minted}{text}
integrate(x * exp(x), x)
integrate(x^a, x)
integrate(sin(x)^2, x)
\end{minted}
  Open the derivation of the first: the by-parts steps are \status{checked}, the check step
  carries the differentiation, and the compare step shows both sides expanded. The second's
  status names the side condition. The third is refused.
\end{trybox}

\section{Exercises}

\begin{exercisebox}[Reading the hypotheses]
  For $\int x^a\,dx$, list the rules that fire in the check's differentiation and say which
  hypothesis of \lc{integrate_deriv} each one's status discharges.
\end{exercisebox}

\begin{exercisebox}[A checker for something else]
  Design a checked command for solving $ax + b = 0$: what does the finder guess, what does the
  check normalize, and what is the analogue of \lc{cmdIntegrate_spec}?
\end{exercisebox}

\chapter{Two More Worlds}\label{ch:worlds}

The engine's term type, its rewriter, its printer, its origin tracker and its notebook were built
for algebra. The last week of the project asked whether they were built for anything else, by
adding two small worlds from the author's other books: the untyped $\lambda$-calculus, and finite
partial orders with joins, meets and fixed points. The answer is yes, with one budget and a
different shape of theorem.

\section{The lambda calculus}

A cell is a $\lambda$-cell if it contains a $\lambda$ or a backslash, a \texttt{:=}, or starts with
a defined name; the engine decides. The parser reads $\lambda x\,y.\,e$, juxtaposition, digits as
Church numerals, and \texttt{name := term}; the Church library --- \texttt{succ}, \texttt{add},
\texttt{mul}, \texttt{true}, \texttt{false}, \texttt{if}, pairs --- is preloaded.

Reduction is normal order, one $\beta$-step at a time, each step recorded:

\begin{minted}{lean}
def betaStep : Term → Option (Term × Bool)
  | .app (.lam x body) arg => some (subst x arg body)
  | .app a b =>
    match betaStep a with
    | some (a', r) => some (.app a' b, r)
    | none => (betaStep b).map fun (b', r) => (.app a b', r)
  | .lam x e => (betaStep e).map fun (e', r) => (.lam x e', r)
  | .var _ => none
\end{minted}

The Boolean says whether an $\alpha$-renaming preceded the step, so that the notebook can label
it \rn{lambda.alpha-beta} rather than \rn{lambda.beta}.

\paragraph{Substitution had to be total.} The textbook capture-avoiding substitution renames a
binder and recurses on the \emph{renamed} body, and Lean cannot see that terminate --- the
author's earlier book marks it \lc{partial}. The engine carries the renaming down in one
structural pass first (\lc{freshen}, with the clash set and the renaming as parameters) and then
substitutes without ever renaming. It computes the same thing, and it is total.

\paragraph{The one budget.} Whether a term has a normal form is undecidable, so reduction runs on
fuel: $\Omega$ is refused after 1000 $\beta$-steps with what it had become. It is the only budget
in the engine, and it is the honest one; everything in Part~III was about removing the others.

\paragraph{Encoded, not extended.} Terms are encoded into \lc{Expr} --- \lc{fn "λ" [var x, body]}
and \lc{fn "@" [f, a]} --- so selection, explanation and origin tracking work unchanged and the
printer learned two heads. The de Bruijn view is computed alongside every step and sent with it;
the notebook's View menu toggles a rendering the engine already delivered. A normal form that is a
Church numeral or a Boolean is read out beside the result, and the reader is proved right on the
engine's own numerals:

\begin{minted}{lean}
theorem readChurch_church (n : Nat) : readChurch (church n) = some n
\end{minted}

What is not proved --- that the renaming preserves $\alpha$-equivalence, that the de Bruijn view
commutes with $\beta$ --- is why the $\beta$-steps are reported \status{unverified}. The
substitution lemma that \emph{is} proved says only that substitution introduces no new free
variables. Appendix~\ref{app:rules} lists all of it.

\section{Finite partial orders}

The order world is a different shape of engine: decisions over lists, and theorems that say each
decision means the textbook Prop. A poset is given as elements and a relation,
\texttt{poset(\{a,b,c\}; a<b, a<c)}; the engine takes the reflexive-transitive closure and decides
reflexivity, antisymmetry and transitivity, naming a witness on failure. Then: covers (the Hasse
diagram), upper and lower bounds, join and meet, lattice-or-witness-pair, top and bottom, maximal
and minimal, \texttt{le} explained by a chain of covers; maps as tables, \texttt{monotone} with a
witness; \texttt{lfp} and \texttt{gfp} by iterating from $\bot$ or $\top$, with the Kleene chain as
the derivation. \texttt{divisors(n)}, \texttt{subsets(\{…\})} and \texttt{chain(n)} build the
standard examples.

The theorems, all Init-only:

\begin{minted}{lean}
theorem checkPartialOrder_none {P : Poset} (h : checkPartialOrder P = none) :
    (∀ x ∈ P.elems, P.rel x x = true) ∧
    (∀ p ∈ P.le, p.1 ≠ p.2 → P.rel p.2 p.1 = false) ∧
    (∀ p ∈ P.le, ∀ q ∈ P.le, p.2 = q.1 → P.rel p.1 q.2 = true)

theorem covers_spec (P : Poset) (x y : String) :
    covers P x y = true ↔ P.lt x y = true ∧ ∀ z ∈ P.elems, ¬ (P.lt x z = true ∧ P.lt z y = true)

theorem sup_spec {P : Poset} {xs : List String} {j : String} (h : sup P xs = some j) :
    j ∈ upperBounds P xs ∧ ∀ v ∈ upperBounds P xs, P.rel j v = true

theorem iter_le_fixed {P : Poset} {f : PMap} (hm : monotoneFailure P f = none) {y : String}
    (hy : f.apply y = y) : ∀ (n : Nat) (x : String), P.rel x y = true → P.rel (f.iter n x) y = true
\end{minted}

The last is the Knaster--Tarski fragment a finite poset needs: every point of the Kleene chain
lies below every fixed point above its start, so with $x = \bot$ the chain's last element --- a
fixed point, as checked --- is the least. The honest boundary is that the chain stabilizes within
$|P|$ steps; that is pigeonhole on a finite chain, and it is checked at run time rather than
proved.

\begin{logbox}[The same shape as \lcn{checked}]
  Every theorem in this world is ``the decision means the Prop''. It is the shape of the
  run-time boundaries of Chapter~\ref{ch:order}, but here the Prop is the definition a reader
  meets in a first course, which is what should sit next to a Hasse diagram. The notebook draws
  the diagram in layers by height from a list of covers, and \lc{covers_spec} is the reason the
  picture is right.
\end{logbox}

\begin{trybox}
\begin{minted}{text}
let D = divisors(12)
join(D, 4, 6)
lattice(D)
let f = map(D; 1->2, 2->2, 3->6, 4->4, 6->6, 12->12)
lfp(D, f)
add 2 3
(λx. x x) (λx. x x)
\end{minted}
  The Kleene chain is the derivation of \texttt{lfp}. Type \texttt{\textbackslash lam} and
  Tab for the $\lambda$. The result of \texttt{add 2 3} reads ``the Church numeral 5''; toggle
  de Bruijn indices in the View menu. The last line is $\Omega$, refused after 1000 steps with
  what it had become.
\end{trybox}

\section{Exercises}

\begin{exercisebox}[The renaming pass]
  Show that \lc{freshen} followed by \lc{substRaw} computes the same term as the textbook's
  recursive-renaming substitution on $(\lambda y.\, x\, y)[x := y]$, and say what the textbook
  version recurses on that Lean cannot see decrease.
\end{exercisebox}

\begin{exercisebox}[Pigeonhole, proved]
  State the theorem that closes the order world's boundary: on a finite poset, a monotone map
  iterated from $\bot$ reaches a fixed point within $|P|$ steps. Which lemma about
  \lc{List.Nodup} would you reach for?
\end{exercisebox}

\part{The Ledger}
\chapter{What Is Proved, and What Is Not}\label{ch:ledger}

The book has used four words throughout --- verified, conditional, checked, unverified --- and this
chapter is their ledger. It is also the argument of the book stated once, plainly: a verified
system is not one where everything is proved, but one where the boundary between proved and not is
drawn in the code, reported to the user, and never crossed silently.

\section{The four statuses, once more}

\begin{itemize}[nosep]
  \item \status{verified}: an unconditional theorem, over $\mathbb{R}$ where the rule has a real
    meaning, over $\mathbb{Q}$ for elimination, over the Prop for the order world.
  \item \status{conditional}: a theorem with a side condition whose necessity is itself proved by
    a counterexample. The notebook shows the condition.
  \item \status{checked}: a guess that a later step verifies, and the later step is verified.
  \item \status{unverified}: no theorem, and the note says what would be needed.
\end{itemize}

Appendix~\ref{app:rules} is the table, generated from the same list the engine sends the notebook,
so the book and the software cannot disagree about it.

\section{Termination}

Cell evaluation has no step budget. Every rule of the notebook pipeline decreases the six-tier
ordering on a node whose children are normal (\lc{pipelineOrdered}), for every checker the
pipeline is instantiated with (\lc{pipelineOrderedWith}). Expansion is a structurally recursive
function. Elimination is a loop over columns. The one remaining budget is $\beta$-reduction, whose
termination is undecidable, and it is refused loudly at 1000 steps.

Where the proof needed a property of a subsystem's output it does not see inside, that property is
decided at run time and a failure is an error: a command's output has no command in it; a radical
rule's guard holds; the Kleene chain stabilized. These are boundaries, not budgets. A boundary is a
property of a result; a budget is a count of steps; the difference is whether the failure can be
explained.

\section{Axioms}

Every theorem in \file{engine/} and \file{proofs/} depends on the three standard axioms of Lean's
classical mathematics and no others: propositional extensionality, choice, and quotient soundness.
There are no \lc{sorry}s. \lc{#print axioms} on the top-level theorems is part of the build.

\section{Sizes}

\begin{center}
\begin{tabular}{lr}
  Lean source (engine and proofs) & 11\,389 lines \\
  Theorems and lemmas & 547 \\
  WebAssembly module & 1.8 MB \\
  First reply in the browser & 96 ms \\
\end{tabular}
\end{center}

\section{What is not proved}

The list is in the table, but the pattern is worth stating. What is unproved is either
\emph{vacuous} (a theorem about matrices in a semantics where matrices have no value), an
\emph{abbreviation} (nothing to prove), \emph{symbolic} where the verified path is numeric, or
\emph{metatheory} that the author has not yet written ($\alpha$-equivalence, de Bruijn
commutation, pigeonhole). None of it is a rule that is believed wrong. And every item was on the
open list before it was in this chapter.

\section{What the book argued}

That a learning notebook can show its work and mean it. That the meaning is a theorem per rule,
folded once over the rewriter and instantiated per semantics. That termination without a budget
is an ordering, that the ordering is a design, and that what the ordering cannot afford is better
done as a function than faked as a rule. That a guess plus a verified check is a verified
answer. And that the honest label on a step is worth more than a proof the user cannot see.

The notebook is at the address on the title page. Every ``Try it'' box in this book runs there,
in the browser, on the engine whose theorems this chapter counted.


\appendix
\chapter{The Rule Table}\label{app:rules}

The statuses below are the ones the engine reports through \texttt{engine.capabilities}. A rule
the engine omits is unverified.

\newcommand{\rt}[3]{\rn{#1} & \status{#2} & #3 \\}
\footnotesize
\begin{longtable}{p{4.3cm}p{2.15cm}p{7.65cm}}
\toprule Rule & Status & Note \\ \midrule \endhead
\multicolumn{3}{l}{\textbf{Simplification}} \\
\rt{simp.flatten}{verified}{Associativity of $+$ and $\cdot$.}
\rt{simp.identity}{verified}{The additive and multiplicative identities and the annihilator.}
\rt{simp.fold-constants}{verified}{Exact rational arithmetic; $\mathbb{Q}$ embeds in $\mathbb{R}$.}
\rt{simp.collect-like-terms}{verified}{Distributivity: $a t + b t = (a+b) t$.}
\rt{simp.power}{verified}{Includes exact roots; the root search returns only checked roots.}
\rt{simp.collect-powers}{conditional}{$b^m b^n = b^{m+n}$ needs a positive base: at $b = 0$ it turns $0$ into $1$.}
\rt{simp.function}{conditional}{$\exp(\ln x) = x$ needs $0 < x$: at $x = -1$ it turns $-1$ into $1$. The other cases, $\sin/\cos = \tan$ among them, are unconditional.}
\rt{simp.radical}{verified}{A perfect-power base is reduced ($8^{1/2} = 2^{3/2}$) and same-index radicals multiply under one root; unconditional, the bases being positive integers.}
\rt{simp.collect-radicals}{verified}{Radicals with the same square-free part collect, $\sqrt{50} - \sqrt{18} = 2\sqrt 2$; unconditional.}
\midrule \multicolumn{3}{l}{\textbf{Derivatives}} \\
\rt{diff.constant}{verified}{A term the variable does not occur in has derivative $0$.}
\rt{diff.variable}{verified}{The identity function has slope $1$ everywhere.}
\rt{diff.constant-multiple}{verified}{Constant factors pull out; no differentiability needed.}
\rt{diff.sum}{conditional}{Needs every summand differentiable: for $|x| + x$ at $0$ the rule gives $1$ where the derivative does not exist.}
\rt{diff.product}{conditional}{Needs both factors differentiable at the point.}
\rt{diff.power}{conditional}{Proved for a natural exponent with a differentiable base; real exponents still open.}
\rt{diff.chain}{conditional}{Needs the inner function differentiable; proved for $\sin$, $\cos$ and $\exp$, while $\ln$ and $\tan$ also need a domain condition.}
\rt{diff.matrix}{unverified}{Proved, but matrices carry no value in the $\mathbb{R}$ semantics, so the theorem has no content yet.}
\rt{diff.higher-order}{unverified}{An abbreviation: it eliminates the three-argument form, so there is nothing to prove.}
\midrule \multicolumn{3}{l}{\textbf{Linear algebra}} \\
\rt{cmd.rref}{verified}{Over $\mathbb{Q}$ the reduced matrix has the input's solution set (\lc{sol_rref}) and is in reduced row echelon form (\lc{rref_isRref}). With symbolic entries the nested row operations are the unverified \texttt{.symbolic} ones, and the step inherits their status.}
\rt{la.row-swap}{verified}{Exchanging two rows preserves the solution set (\lc{sol_swap}).}
\rt{la.row-scale}{verified}{Scaling a row by a nonzero rational preserves the solution set (\lc{sol_scale}).}
\rt{la.row-add}{verified}{Adding a multiple of another row preserves the solution set (\lc{sol_addMul}).}
\rt{la.row-swap.symbolic}{unverified}{Symbolic entries: the pivot is assumed nonzero because simplify could not show it is zero.}
\rt{la.row-scale.symbolic}{unverified}{Symbolic entries: division by a pivot that is only assumed nonzero.}
\rt{la.row-add.symbolic}{unverified}{Symbolic entries: the arithmetic is the simplifier's, outside the $\mathbb{Q}$ theorem.}
\midrule \multicolumn{3}{l}{\textbf{Expansion}} \\
\rt{cmd.expand}{verified}{Distribution is a total function proved sound over $\mathbb{R}$ (\lc{dist_sound}); the collection afterwards is the pipeline's own steps with their statuses.}
\rt{expand.distribute}{verified}{Multiplying out a product of sums and collecting like monomials: \lc{dist_sound}.}
\rt{expand.power}{verified}{A power of a sum is the sum multiplied by itself: \lc{dist_sound}.}
\midrule \multicolumn{3}{l}{\textbf{Integration}} \\
\rt{cmd.integrate}{verified}{Accepted only when the candidate's derivative normalizes to the integrand, exactly (\lc{cmdIntegrate_spec}); \lc{integrate_deriv} reads that as $F' = f$ wherever the differentiation steps shown are sound. The finder's own steps are guesses.}
\rt{int.check}{verified}{The differentiation of the candidate: this step carries the claim, with the statuses of its own steps.}
\rt{int.compare}{verified}{Derivative and integrand are expanded (\lc{dist}, proved sound) and simplified before comparison; the statuses of the simplification steps apply.}
\rt{int.constant}{checked}{A guess from the finder; nothing is proved about it. Accepted only because \rn{int.check} verifies the result by differentiation.}
\rt{int.variable}{checked}{A guess from the finder, verified by \rn{int.check}.}
\rt{int.sum}{checked}{A guess from the finder, verified by \rn{int.check}.}
\rt{int.constant-multiple}{checked}{A guess from the finder, verified by \rn{int.check}.}
\rt{int.power}{checked}{A guess from the finder, verified by \rn{int.check} (the symbolic-exponent case relies on \rn{simp.collect-powers} there).}
\rt{int.exponential}{checked}{A guess from the finder, verified by \rn{int.check}.}
\rt{int.table}{checked}{A guess from the finder, verified by \rn{int.check}.}
\rt{int.linear-substitution}{checked}{A guess from the finder, verified by \rn{int.check}.}
\rt{int.substitution}{checked}{A guess from the finder ($u$-substitution), verified by \rn{int.check}.}
\rt{int.by-parts}{checked}{A guess from the finder (integration by parts), verified by \rn{int.check}.}
\midrule \multicolumn{3}{l}{\textbf{Order world}} \\
\rt{order.closure}{verified}{The order is the reflexive-transitive closure, and reflexivity, antisymmetry and transitivity are decided (\lc{checkPartialOrder_none}).}
\rt{order.covers}{verified}{Hasse edges are exactly the covers: $x < y$ with nothing strictly between (\lc{covers_spec}).}
\rt{order.upper-bounds}{verified}{Every element above both, by the decision on the finite order.}
\rt{order.least}{verified}{The element found is an upper bound below every upper bound (\lc{sup_spec}); none exists when the search fails (\lc{sup_none}).}
\rt{order.lower-bounds}{verified}{Every element below both, by the decision on the finite order.}
\rt{order.greatest}{verified}{Dual of the join: the greatest lower bound.}
\rt{order.lattice}{verified}{Every pair searched for a join and a meet; the witness is reported when one is missing.}
\rt{order.cover}{verified}{A cover in the Hasse diagram; the chain composes by transitivity.}
\rt{order.monotone}{verified}{Every pair $x \le y$ of the order checked; the witness is reported when it fails.}
\rt{order.iterate}{verified}{One step of the Kleene chain; each element of the chain is below every fixed point (\lc{iter_le_fixed}).}
\rt{order.fixed}{verified}{The chain stopped at a fixed point (checked), which \lc{iter_le_fixed} makes the least (dually, the greatest).}
\midrule \multicolumn{3}{l}{\textbf{$\lambda$-calculus}} \\
\rt{lambda.delta}{verified}{Unfolding a definition replaces a free name by its term; nothing to prove beyond that.}
\rt{lambda.beta}{unverified}{$\beta$-reduction with capture-avoiding substitution; the substitution lemma is not yet proved.}
\rt{lambda.alpha-beta}{unverified}{A binder renamed to avoid capture, then $\beta$; the renaming is not yet proved to preserve $\alpha$-equivalence.}
\bottomrule
\end{longtable}
\normalsize

\chapter{Building It}\label{app:build}

The repository is \file{github.com/adekau/mathbook}. The toolchain is the latest stable Lean 4
(\file{engine/lean-toolchain}, v4.33.1 at this commit); \file{proofs/} pins Mathlib at the
release tag whose toolchain matches, and \file{engine/} depends on \lc{Init} alone, which the
script \file{scripts/check-engine-deps.sh} enforces in CI.

\section*{The engine and its tests}
\begin{minted}{text}
cd engine && lake build && lake test
\end{minted}
\texttt{lake test} replays \file{Tests/golden.tsv}, the wire-level answers recorded by the
differential test against the original TypeScript engine before that engine was deleted, and
the unit tests in \file{Tests/Main.lean}.

\section*{The proofs}
\begin{minted}{text}
cd proofs && lake exe cache get && lake build
\end{minted}
The first command fetches Mathlib's compiled cache; the build then takes a few minutes. Every
top-level theorem's axioms are printed as part of it.

\section*{The notebook}
\begin{minted}{text}
npm install && npm test && npm run bundle
\end{minted}
\texttt{npm test} runs the TypeScript tests, including the native stdio round trip against the
Lean executable. \texttt{npm run bundle} writes the static notebook to
\file{apps/notebook/dist}, with a build stamp on every asset.

\section*{WebAssembly}
\begin{minted}{text}
npm run wasm
\end{minted}
Builds the Lean runtime and \lc{Init} for \texttt{wasm32} from the toolchain's sources (once;
cached under \file{engine/toolchains/}), then links the engine. The first build takes about
ten minutes, most of it elaborating \lc{Init}; the patch the runtime needs is in
\file{engine/wasm/}.

\section*{This book}
\begin{minted}{text}
cd book && latexmk -xelatex -shell-escape show-your-work.tex
\end{minted}
XeLaTeX with \texttt{minted} (Pygments) and \texttt{tcolorbox}. The CI workflow
\file{.github/workflows/build-book.yml} builds it on every push and attaches the PDF to a
release tagged with the date and commit.


\end{document}
